\documentclass[10pt,a4paper]{amsart}
\usepackage{amsmath,amssymb,amsthm,graphicx}
% Retained preamble for the Stacks Project chapter snapshot in this directory.
% The chapter's own preamble.tex is a copy of the shared Stacks Project preamble; the block
% below is sliced MECHANICALLY, line for line, from it. Nothing was added, removed or
% redefined; the wrapper class, the xr-hyper/hyperref wiring and the externaldocument list
% are omitted because the assembler supplies the document class.
%
% lines taken from preamble.tex: 3, 5, 9, 11, 12, 156, 157, 158, ... (51 lines total)
%
% package coverage restored from the chapter's own preamble (2026-10-13): xy, xy-2cell

\usepackage[all]{xy}
\xyoption{2cell}
\UseAllTwocells

\newenvironment{reference}{\comment}{\endcomment}
\newenvironment{slogan}{\comment}{\endcomment}
\newenvironment{history}{\comment}{\endcomment}
\theoremstyle{plain}
\newtheorem{theorem}[subsection]{Theorem}
\newtheorem{proposition}[subsection]{Proposition}
\newtheorem{lemma}[subsection]{Lemma}
\theoremstyle{definition}
\newtheorem{definition}[subsection]{Definition}
\newtheorem{example}[subsection]{Example}
\newtheorem{exercise}[subsection]{Exercise}
\newtheorem{situation}[subsection]{Situation}
\theoremstyle{remark}
\newtheorem{remark}[subsection]{Remark}
\newtheorem{remarks}[subsection]{Remarks}
\numberwithin{equation}{subsection}
\def\lim{\mathop{\mathrm{lim}}\nolimits}
\def\colim{\mathop{\mathrm{colim}}\nolimits}
\def\Spec{\mathop{\mathrm{Spec}}}
\def\Hom{\mathop{\mathrm{Hom}}\nolimits}
\def\Ext{\mathop{\mathrm{Ext}}\nolimits}
\def\SheafHom{\mathop{\mathcal{H}\!\mathit{om}}\nolimits}
\def\SheafExt{\mathop{\mathcal{E}\!\mathit{xt}}\nolimits}
\def\Sch{\mathit{Sch}}
\def\Mor{\mathop{\mathrm{Mor}}\nolimits}
\def\Ob{\mathop{\mathrm{Ob}}\nolimits}
\def\Sh{\mathop{\mathit{Sh}}\nolimits}
\def\NL{\mathop{N\!L}\nolimits}
\def\CH{\mathop{\mathrm{CH}}\nolimits}
\def\proetale{{pro\text{-}\acute{e}tale}}
\def\etale{{\acute{e}tale}}
\def\QCoh{\mathit{QCoh}}
\def\Ker{\mathop{\mathrm{Ker}}}
\def\Im{\mathop{\mathrm{Im}}}
\def\Coker{\mathop{\mathrm{Coker}}}
\def\Coim{\mathop{\mathrm{Coim}}}
\DeclareMathSymbol{\boxtimes}{\mathbin}{AMSa}{"02}
\def\QCohstack{\mathcal{QC}\!\mathit{oh}}
\def\Cohstack{\mathcal{C}\!\mathit{oh}}
\def\Spacesstack{\mathcal{S}\!\mathit{paces}}
\def\Quotfunctor{\mathrm{Quot}}
\def\Hilbfunctor{\mathrm{Hilb}}
\def\Curvesstack{\mathcal{C}\!\mathit{urves}}
\def\Polarizedstack{\mathcal{P}\!\mathit{olarized}}
\def\Complexesstack{\mathcal{C}\!\mathit{omplexes}}
\def\Pic{\mathop{\mathrm{Pic}}\nolimits}
\def\Picardstack{\mathcal{P}\!\mathit{ic}}
\def\Picardfunctor{\mathrm{Pic}}
\def\Deformationcategory{\mathcal{D}\!\mathit{ef}}
\usepackage{cleveref}
\crefname{originallemma}{original lemma}{original lemmas}
\Crefname{originallemma}{Original lemma}{Original lemmas}
\makeatletter
\newlabel{algebra-definition-lci-field}{{00S9}{source}{}{sourcepointer-1}{}}
\newlabel{algebra-definition-lci-field@cref}{{[originallemma][0][]00S9}{[1][0][]source}{}{}{}}
\makeatother
\crefname{originallemma}{original lemma}{original lemmas}
\Crefname{originallemma}{Original lemma}{Original lemmas}
\makeatletter
\newlabel{algebra-lemma-grothendieck-regular-sequence}{{00MG}{source}{}{sourcepointer-2}{}}
\newlabel{algebra-lemma-grothendieck-regular-sequence@cref}{{[originallemma][0][]00MG}{[1][0][]source}{}{}{}}
\makeatother
\crefname{originallemma}{original lemma}{original lemmas}
\Crefname{originallemma}{Original lemma}{Original lemmas}
\makeatletter
\newlabel{algebra-lemma-lci}{{00SC}{source}{}{sourcepointer-3}{}}
\newlabel{algebra-lemma-lci@cref}{{[originallemma][0][]00SC}{[1][0][]source}{}{}{}}
\makeatother
\crefname{originallemma}{original lemma}{original lemmas}
\Crefname{originallemma}{Original lemma}{Original lemmas}
\makeatletter
\newlabel{lemma-jouanolou-type-thing}{{0E7U}{source}{}{sourcepointer-4}{}}
\newlabel{lemma-jouanolou-type-thing@cref}{{[originallemma][0][]0E7U}{[1][0][]source}{}{}{}}
\makeatother
\title{Stacks, Examples of Deformations: lemma-smoothing-affine-lci}
\author{The Stacks Project Authors}
\date{stacks-project:examples-defos.tex:local-snapshot-20261012; GNU FDL 1.2 or later}
\begin{document}
\maketitle
\noindent\textit{Editorial note.} Statement span is lines 3215-3221 of sources/stacks-examples-defos/pinned-original-source.tex; proof span is lines 3223-3258, the complete author proof (34 body lines).\par\medskip
\noindent\textit{Editorial note.} Navigation stubs: 4 cited-but-not-included labels are supplied as non-load-bearing pointers with their pinned tags where the snapshot tag index has them.\par\medskip
\noindent\textit{Editorial note.} extra\_packages is empty on purpose: the retained preamble already loads xy with options and a second declaration clashes in the assembler wrapper.\par\medskip
\noindent\textit{Editorial note.} Text is retained verbatim, including the author's own typos and placeholders.\par\medskip
\noindent\textit{Editorial note.} Queue index 16 of queue\_examples-defos.json (pb=1436, stmt 3215-3221, proof 3223-3258).\par\medskip
\noindent\textit{Allowed standard prerequisites (editorial).}\begin{itemize}
\item Explicit prerequisite boundary. The cited inputs are prerequisites; the argument in the retained span is not replaced by any cited result.
\end{itemize}\medskip
\ifdefined\hypertarget\hypertarget{sourcepointer-1}{}\fi
\noindent\textit{Source reference.} Original lemma 00S9: sources/stacks-examples-defos/pinned-original-source.tex, source lines 3258--3258. Cited by the retained statement or proof but not part of the collected spans; supplied only as a navigation pointer. The argument does not depend on its content.\par\medskip
\ifdefined\hypertarget\hypertarget{sourcepointer-2}{}\fi
\noindent\textit{Source reference.} Original lemma 00MG: sources/stacks-examples-defos/pinned-original-source.tex, source lines 3258--3258. Cited by the retained statement or proof but not part of the collected spans; supplied only as a navigation pointer. The argument does not depend on its content.\par\medskip
\ifdefined\hypertarget\hypertarget{sourcepointer-3}{}\fi
\noindent\textit{Source reference.} Original lemma 00SC: sources/stacks-examples-defos/pinned-original-source.tex, source lines 3258--3258. Cited by the retained statement or proof but not part of the collected spans; supplied only as a navigation pointer. The argument does not depend on its content.\par\medskip
\ifdefined\hypertarget\hypertarget{sourcepointer-4}{}\fi
\noindent\textit{Source reference.} Original lemma 0E7U: sources/stacks-examples-defos/pinned-original-source.tex, source lines 3151--3166. Cited by the retained statement or proof but not part of the collected spans; supplied only as a navigation pointer. The argument does not depend on its content.\par\medskip
\begin{lemma}
\label{lemma-smoothing-affine-lci}
Let $k$ be a field. Let $A$ be a global complete intersection
over $k$. There exists a flat finite type ring map
$k[[t]] \to B$ with $B/tB \cong A$ such that
$B[1/t]$ is smooth over $k((t))$.
\end{lemma}

\begin{proof}
Write $A = k[x_1, \ldots, x_n]/(f_1, \ldots, f_c)$ as in
Algebra, Definition \ref{algebra-definition-lci-field}.
We are going to choose
$a_{ij} \in (t) \subset k[[t]]$ and set
$$
g_j = f_j + a_{0j} + a_{1j}x_1 + \ldots + a_{nj}x_n \in
k[[t]][x_1, \ldots, x_n]
$$
After doing this we take
$B = k[[t]][x_1, \ldots, x_n]/(g_1, \ldots, g_c)$.
We claim that $k[[t]] \to B$ is flat at every prime ideal
lying over $(t)$. Namely, the elements $f_1, \ldots, f_c$
form a regular sequence in the local ring at any prime ideal
$\mathfrak p$ of $k[x_1, \ldots, x_n]$ containing $f_1, \ldots, f_c$
(Algebra, Lemma \ref{algebra-lemma-lci}). Thus $g_1, \ldots, g_c$
is locally a lift of a regular sequence and we can apply
Algebra, Lemma \ref{algebra-lemma-grothendieck-regular-sequence}.
Flatness at primes lying over $(0) \subset k[[t]]$ is automatic
because $k((t)) = k[[t]]_{(0)}$ is a field. Thus $B$ is flat
over $k[[t]]$.

\medskip\noindent
All that remains is to show that for suitable choices
of $a_{ij}$ the generic fibre $B_{(0)}$ is smooth over
$k((t))$. For this we have to show that we can choose
our $a_{ij}$ so that the induced morphism
$$
(a_{ij}) : \Spec(k[[t]]) \longrightarrow \mathbf{A}^{c(n + 1)}_k
$$
maps into the nonempty Zariski open of
Lemma \ref{lemma-jouanolou-type-thing}.
This is clear because there is no nonzero polynomial in the
$a_{ij}$ which vanishes on $(t)^{\oplus c(n + 1)}$.
(We leave this as an exercise to the reader.)
\end{proof}

\end{document}
